\documentclass{article}
\usepackage{amssymb,amsmath}
\usepackage{graphicx}
\usepackage{enumerate}
\usepackage{amssymb,amsmath,amsthm}
\usepackage{graphicx}
\usepackage{tikz}
\usepackage{pgfplots}
\usepackage[utf8]{inputenc}
\usepackage[brazil]{babel}
\usepackage{enumerate}
\usepackage{color}
\usepackage{mathabx}
\usepackage{calc}
\usepackage{fullpage}

\usepackage{dsfont} % includes bb 1
\newcommand*\1{\mathds{1}}
\usepackage[brazil]{babel}

\newtheorem{theorem}{Teorema}[section]
\newtheorem{corollary}[theorem]{Corolário}
\newtheorem{lemma}[theorem]{Lema}
\newtheorem{proposition}[theorem]{Proposição}
\newtheorem{definition}[theorem]{Definição}
\newtheorem{notation}[theorem]{Notação}

\DeclareMathOperator{\Var}{Var}
\def\d{\mathrm{d}}

\begin{document}

\title{Probabilidade I}
\author{Segunda Prova}

\maketitle

\noindent Observa\c{c}\~oes:
\begin{itemize}
\item A prova terá duração de $3$ horas e meia.
\item Questões deverão ser respondidas de maneira matematicamente rigorosa e clara.
\item Todos os resultados provados em aula (exceto exercícios) poderão ser utilizados, salvo quando a questão pedir que algum destes seja provado. Nesse caso, todos os resultados anteriores poderão ser utilizados.
\item Caso uma questão seja composta de vários itens, cada um poderá ser resolvido independentemente, supondo a validade dos anteriores.
\end{itemize}

\vspace{4mm}
\noindent {\bf 1)}
Seja $(X,Y)$ um vetor aleatório em $\mathbb R^n\times\mathbb R^m$ com densidade conjunta
\begin{equation}
    \rho(x,y)
\end{equation}
em relação à medida de Lebesgue em $\mathbb R^n\times\mathbb R^m$.

\begin{enumerate}[\quad a)]
\item Mostre que $P_X$ e $P_Y$ possuem densidades
\begin{equation}
  \rho_X(x)=\int_{\mathbb R^m}\rho(x,y)\,dy,
  \qquad
  \rho_Y(y)=\int_{\mathbb R^n}\rho(x,y)\,dx.
\end{equation}
Em particular, para todo boreliano $A\subseteq \mathbb R^n$,
  \begin{equation}
    \mathbb P(X\in A) = \int_A \rho_X(x)\,dx.
  \end{equation}
\item Defina
  \begin{equation}
    \rho_{X \mid Y}(x, y)
    = \begin{cases}
        \frac{\rho(x,y)}{\rho_Y(y)} & \qquad \text{se $\rho_Y(y) > 0$.}\\
        \rho_X(x) & \text{caso contrário}
      \end{cases}.
    \end{equation}
Mostre que $\rho_{X \mid Y}(x, y) dx$ é a probabilidade condicional regular $P[X \in \cdot | Y = y]$.

    \item Deduza a fórmula de Bayes para densidades:
    \begin{equation}
    \rho_{X \mid Y}(x, y)
    =
    \frac{\rho_{Y \mid X}(x, y)\rho_X(x)}
    {\int_{\mathbb R^n}\rho_{Y\mid X}(z, y)\rho_X(z) \, dz}.
    \end{equation}
\end{enumerate}

\pagebreak

\noindent {\bf 2)} Seja $(\lambda_i)_{i \ge 0}$ uma sequência de números positivos tal que
\begin{equation}
\sup_{m \ge 0} \lambda_m = \Lambda < \infty.
\end{equation}

Considere o seguinte candidato a processo de Markov em tempo contínuo $(X_t)_{t \ge 0}$ em $\mathbb N$ que, quando está no estado $m$, salta para $m + 1$ com taxa $\lambda_m$.

Fixado um inteiro $n \geq 0$, vamos definir um espaço de probabilidade denotado por $\mathbb{P}_n$.
Para isso, sejam $E_0, E_1, E_2, \dots$ variáveis aleatórias independentes tais que
\begin{equation}
  E_k \sim \mathrm{Exp}(\lambda_{n + k}), \text{ para $k \geq 0$}.
\end{equation}
Defina
\begin{equation}
  T_0 = 0,
  \qquad
  T_k = E_0 + \cdots + E_{k - 1}, \quad \text{ para $k \ge 1$}
\end{equation}
e
\begin{equation}
  X_t = n + k
  \quad
  \text{se}
  \quad
  T_k\le t<T_{k+1}.
\end{equation}

\begin{enumerate}[\quad a)]
\item Mostre que, sob $\mathbb{P}_n$, $(X_t)_{t \ge 0}$ está bem definido para todo $t \ge 0$, isto é, que
  \begin{equation}
    T_k\to+\infty
    \quad
    \text{q.c.}
  \end{equation}
  Sugestão: use que $\lambda_n\le \Lambda$ para todo $n$.

\item Mostre que, sob $\mathbb{P}_n$, para qualquer $m \geq n$, condicionalmente a $X_t = m$, o tempo restante até o próximo salto tem distribuição exponencial de parâmetro $\lambda_m$.

\item Fixado $n \geq 0$, gostaríamos de concluir que, sob a lei $\mathbb{P}_0$ condicionada a $X_t = n$, o processo futuro
  \begin{equation}
    (X_{t + s})_{s \ge 0}
  \end{equation}
  tem a mesma lei que o processo iniciado em $n$.
  Mas por simplicidade, mostre que $(X_t)_{t\ge 0}$ satisfaz uma versão bastante fraca da propriedade de Markov:
  \begin{equation}
    \mathbb P_0(X_{t + s} = m \mid X_t = n)
    =
    \mathbb P_n(X_s = m),
    \qquad s, t \ge 0.
  \end{equation}
\end{enumerate}

\vspace{4mm}
\noindent {\bf 3)} Sejam $X_1,X_2,\dots$ variáveis aleatórias iid tais que
\begin{equation}
\mathbb P(X_1=1)=p,
\qquad
\mathbb P(X_1=-1)=q,
\end{equation}
onde $p<q$ e $p+q=1$.

Defina
\begin{equation}
S_n=X_1+\cdots+X_n,
\qquad
\mathcal F_n=\sigma(X_1,\dots,X_n).
\end{equation}
Dizemos que uma família de variáveis aleatórias $Y_n \in \mathcal{L}^1$ é um martingal com respeito a $\mathcal{F}_n$ se
\begin{equation}
  \mathbb{E}[Y_{n + 1} | \mathcal{F}_n] = Y_n, \quad \text{para todo $n \geq 0$}.
\end{equation}

\begin{enumerate}
    \item Mostre que
    \begin{equation}
    M_n=S_n-(p-q)n
    \end{equation}
    é uma martingal.

    \item Encontre uma constante $r>1$ tal que
    \begin{equation}
    Z_n=r^{S_n}
    \end{equation}
    seja uma martingal.
\end{enumerate}
\end{document}
